\chapter{Access Control Lists}\label{ch:acl}
\bp{4.6}

\begin{outcomes}
By the end of this chapter you will be able to:
\begin{tight}
  \item \textbf{Calculate} a wildcard mask for any subnet, and recognise the
        common ones on sight.
  \item \textbf{Configure} standard, extended, numbered and named IPv4 ACLs.
  \item \textbf{Decide} where to place an ACL and in which direction, and justify
        the choice.
  \item \textbf{Read} an existing ACL and predict exactly which traffic it permits
        --- which is what the exam asks far more often than it asks you to write
        one.
\end{tight}
\end{outcomes}

\begin{prereq}
\chref{ipv4} for subnet masks --- a wildcard mask is meaningless without them.
\chref{transport} for protocol numbers and well-known ports. \chref{nat} for the
order of operations, which decides which address an ACL actually sees.
\end{prereq}

\begin{keyterms}
access control list $\bullet$ wildcard mask $\bullet$ implicit deny $\bullet$
standard ACL $\bullet$ extended ACL $\bullet$ named ACL $\bullet$ sequence number
$\bullet$ \cmd{ip access-group} $\bullet$ \cmd{access-class} $\bullet$
\cmd{established}
\end{keyterms}

\begin{examnote}[IPv4 only, and all four kinds]
Topic~\tp{4.6} reads \emph{configure IPv4 access control lists (standard,
extended, numbered, and named ACLs)}. IPv6 ACLs are not on this blueprint.

All four named kinds are examinable, but note they are two axes, not four
categories: an ACL is standard \emph{or} extended by what it can match, and
numbered \emph{or} named by how you refer to it. All four combinations exist.
\end{examnote}

\section{What an ACL does}

\begin{definitionbox}[Access control list]
An ordered list of \cmd{permit} and \cmd{deny} statements. Each packet is tested
against the entries from the top down; the \textbf{first match wins} and testing
stops. A packet matching nothing hits an invisible \cmd{deny any} at the end and
is dropped.
\end{definitionbox}

Three consequences follow from that definition, and between them they cause most
ACL faults.

\begin{tight}
  \item \textbf{Order is behaviour.} Moving one line changes what the ACL does. A
        specific rule placed after a general one that already matched is dead
        configuration.
  \item \textbf{The implicit deny is always there.} An ACL that permits one thing
        denies everything else. There is no way to switch it off; if you want
        everything else through, end the list with \cmd{permit ip any any}.
  \item \textbf{An unapplied ACL does nothing.} Writing the list and applying it
        to an interface are two separate acts, and forgetting the second is
        common.
\end{tight}

\section{Wildcard masks}

A wildcard mask says which bits of an address must match. A \textbf{0} bit means
``must match''; a \textbf{1} bit means ``ignore''. That is the inverse of a subnet
mask, and the arithmetic is: subtract each octet of the subnet mask from 255.

\begin{longtable}{@{}L{4.2cm}L{3.6cm}L{7.5cm}@{}}
\toprule
\rowcolor{primary!15}
\textbf{Subnet mask} & \textbf{Wildcard} & \textbf{Matches} \\
\midrule
\endhead
\cmd{255.255.255.255} (/32) & \cmd{0.0.0.0} & Exactly one host. Shorthand:
\cmd{host 10.1.1.5} \\
\rowcolor{lightbg}
\cmd{255.255.255.192} (/26) & \cmd{0.0.0.63} & One 64-address subnet \\
\cmd{255.255.255.0} (/24) & \cmd{0.0.0.255} & One 256-address subnet \\
\rowcolor{lightbg}
\cmd{255.255.252.0} (/22) & \cmd{0.0.3.255} & Four consecutive /24s \\
\cmd{255.255.0.0} (/16) & \cmd{0.0.255.255} & An entire class-B-sized range \\
\rowcolor{lightbg}
--- & \cmd{255.255.255.255} & Anything. Shorthand: \cmd{any} \\
\bottomrule
\end{longtable}

\begin{worked}[Wildcard for 172.16.8.0/21]
The mask for a \cmd{/21} is \cmd{255.255.248.0}.

Subtract each octet from 255: $255-255=0$, $255-255=0$, $255-248=7$,
$255-0=255$.

The wildcard is \cmd{0.0.7.255}, and the ACL entry is
\cmd{permit ip 172.16.8.0 0.0.7.255 any}. It matches \cmd{172.16.8.0} through
\cmd{172.16.15.255} --- eight consecutive /24s, which is what a /21 is.
\end{worked}

\begin{conceptbox}[The check that catches arithmetic slips]
Add the wildcard to the network address. The result should be the last address in
the range you meant. $172.16.8.0 + 0.0.7.255 = 172.16.15.255$. If that is not the
broadcast address of the block you intended, the wildcard is wrong.

Do this every time. It takes five seconds and it catches the off-by-one that
otherwise costs you a lab.
\end{conceptbox}

\section{Standard ACLs}

A standard ACL matches the \textbf{source address only}. Numbered 1--99 and
1300--1999.

\begin{config}{a numbered standard ACL}
! Order matters: the specific deny must precede the general permit,
! or the permit matches first and the deny is never reached.
access-list 10 remark Block the lab subnet, allow the rest of campus
access-list 10 deny   192.168.30.0 0.0.0.255
access-list 10 permit 192.168.0.0 0.0.255.255
!
interface GigabitEthernet0/0/1
 ip access-group 10 out
\end{config}

\begin{alertbox}[Place a standard ACL close to the \emph{destination}]
Because it can only see the source address, a standard ACL cannot distinguish
``traffic from this subnet to the finance server'' from ``traffic from this
subnet to anywhere''. Applied near the source, it blocks the host from the entire
network, not just from the thing you wanted protected.

So the rule is: standard ACLs go as close to the destination as possible;
extended ACLs go as close to the source as possible, because they can be precise
enough to be safe there and it saves carrying doomed traffic across the network.
\end{alertbox}

\ccnafig{%
\begin{tikzpicture}[font=\scriptsize]
  \tikzset{
    dv/.style={rounded corners=3pt, draw=#1, fill=#1!8, text=#1!45!black,
               line width=0.9pt, minimum width=1.35cm, minimum height=0.95cm,
               align=center, font=\tiny\bfseries},
    w/.style={line width=0.85pt, neutral!65},
    hdr/.style={font=\scriptsize\bfseries, text=neutral!45!black}}

  % Every prose node gets an explicit text width. Without one the node is as
  % wide as its text, which both collides with its neighbours and drags the
  % picture's bounding box out sideways.
  \node[hdr, text width=13.5cm, align=center] at (5.3,3.5)
      {VLAN 30 must not reach the finance server --- but must still reach
       everything else};

  \node[dv=neutral] (h)   at (0,0.7)    {\faIcon{desktop}\\VLAN 30};
  \node[dv=dom3]    (r1)  at (3.0,0.7)  {\faIcon{route}\\R1};
  \node[dv=dom3]    (r2)  at (7.0,0.7)  {\faIcon{route}\\R2};
  \node[dv=dom1]    (fin) at (10.4,1.75){\faIcon{server}\\Finance};
  \node[dv=dom1]    (oth) at (10.4,-0.4){\faIcon{server}\\Everything\\else};
  \draw[w] (h) -- (r1);  \draw[w] (r1) -- (r2);
  \draw[w] (r2) -- (fin);  \draw[w] (r2) -- (oth);

  % wrong placement -- marker on the link, caption well below it
  \draw[accent!85, line width=1.4pt] (4.7,1.25) -- (5.3,0.15);
  \node[font=\tiny\bfseries, text=accent!75!black, align=center,
        text width=3.1cm, anchor=north] at (5.0,-0.35)
      {\faIcon{times-circle}~\textbf{standard} here\\blocks \emph{everything}};

  % right placement -- caption above, clear of the header
  \draw[dom2!80, line width=1.4pt] (8.5,2.0) -- (9.1,1.1);
  \node[font=\tiny\bfseries, text=dom2!45!black, align=center,
        text width=3.4cm, anchor=south] at (8.8,2.15)
      {\faIcon{check-circle}~\textbf{standard} here\\blocks only this server};

  \node[font=\tiny, text=neutral!85!black, align=center, text width=13.5cm]
      at (5.3,-1.85)
      {A standard ACL sees only \emph{who sent it}, so near the source it cannot
       tell the two servers apart. An \textbf{extended} ACL names the destination
       too, so it is safe at the source --- and saves carrying doomed traffic
       across the network.};
\end{tikzpicture}}%
{Why placement is part of the design. The same ACL is correct in one place and an
outage in the other.}{aclplace}

\section{Extended ACLs}

An extended ACL matches source, destination, protocol and port. Numbered 100--199
and 2000--2699.

The argument order is fixed and worth reciting until it is automatic:

\begin{center}
\cmd{access-list <n> \{permit|deny\} <protocol> <source> <src-port> <dest> <dest-port>}
\end{center}

\begin{config}{an extended ACL letting a subnet reach only what it needs}
access-list 110 remark Guest VLAN -- web, DNS and nothing else
!
! DNS to the campus resolvers, both transports.
access-list 110 permit udp 192.168.40.0 0.0.0.255 host 192.168.20.10 eq 53
access-list 110 permit tcp 192.168.40.0 0.0.0.255 host 192.168.20.10 eq 53
!
! Web out to anywhere.
access-list 110 permit tcp 192.168.40.0 0.0.0.255 any eq 80
access-list 110 permit tcp 192.168.40.0 0.0.0.255 any eq 443
!
! Let guests get an address. Without this the implicit deny stops
! DHCP and the VLAN never works at all.
access-list 110 permit udp any host 255.255.255.255 eq 67
!
! Ping to the gateway only, so users can prove their own link.
access-list 110 permit icmp 192.168.40.0 0.0.0.255 host 192.168.40.1
!
! Everything else is dropped by the implicit deny, logged so that
! surprises are visible.
access-list 110 deny ip any any log
!
interface GigabitEthernet0/0/2
 description Guest VLAN
 ip access-group 110 in
\end{config}

\begin{conceptbox}[\cmd{eq 53} on the destination, not the source]
In that list the port qualifier follows the \emph{destination} address, because
the guest is connecting \emph{to} port 53 on the resolver. Its own source port is
ephemeral and unpredictable, so it is left unspecified.

Getting this backwards --- \cmd{permit udp 192.168.40.0 0.0.0.255 eq 53 any} ---
is a classic and produces an ACL that matches almost nothing. The position of the
port in the line is the whole meaning.
\end{conceptbox}

\begin{examnote}[\cmd{established}]
\cmd{permit tcp any 192.168.10.0 0.0.0.255 established} matches only TCP segments
with the ACK or RST flag set --- that is, packets belonging to a conversation
that somebody inside already started. A pure connection attempt from outside has
only SYN set and is not matched.

It is a cheap way to say ``replies to our traffic may return, new connections may
not.'' It is not stateful --- it inspects flags, not connections --- and the exam
likes that distinction.
\end{examnote}

\section{Named ACLs, and why they are what you should write}

\begin{config}{the same policy, named and editable}
ip access-list extended GUEST-OUT
 remark Guest VLAN -- web, DNS and nothing else
 10 permit udp 192.168.40.0 0.0.0.255 host 192.168.20.10 eq 53
 20 permit tcp 192.168.40.0 0.0.0.255 host 192.168.20.10 eq 53
 30 permit tcp 192.168.40.0 0.0.0.255 any eq 80
 40 permit tcp 192.168.40.0 0.0.0.255 any eq 443
 50 deny   ip any any log
!
interface GigabitEthernet0/0/2
 ip access-group GUEST-OUT in
\end{config}

\begin{alertbox}[This is the practical advantage, and it is a large one]
With sequence numbers you can insert and delete single lines:

\begin{tight}
  \item \cmd{ip access-list extended GUEST-OUT} then \cmd{35 permit tcp ... eq 8080}
        inserts a rule between 30 and 40.
  \item \cmd{no 40} removes one entry and leaves the rest standing.
\end{tight}

With a numbered ACL on older IOS, \cmd{no access-list 110} deletes the
\emph{entire list} --- and if it is still applied to an interface, the interface
now has an empty ACL, which denies everything. That is how a routine edit becomes
an outage. Number your entries in tens so there is room to insert later, and
prefer named ACLs.
\end{alertbox}

\section{Applying and verifying}

\begin{longtable}{@{}L{4.6cm}L{10.5cm}@{}}
\toprule
\rowcolor{primary!15}
\textbf{Command} & \textbf{Where and what} \\
\midrule
\endhead
\cmd{ip access-group <acl> in} &
On an interface. Filters packets \emph{arriving}. Evaluated before routing, so
it sees traffic the router might have dropped anyway --- cheaper. \\
\rowcolor{lightbg}
\cmd{ip access-group <acl> out} &
On an interface. Filters packets \emph{leaving} after a routing decision. Does
not filter traffic the router itself originates. \\
\cmd{access-class <acl> in} &
On \cmd{line vty}. Controls who may open a management session. Use a standard
ACL naming the permitted management subnets. \\
\bottomrule
\end{longtable}

\begin{config}{restricting management access}
ip access-list standard MGMT-HOSTS
 10 permit 10.0.0.0 0.0.0.255
 20 permit host 192.168.99.14
!
line vty 0 15
 access-class MGMT-HOSTS in
 transport input ssh
\end{config}

\begin{verify}{R1}{show access-lists}
Standard IP access list MGMT-HOSTS
    10 permit 10.0.0.0, wildcard bits 0.0.0.255 (214 matches)
    20 permit 192.168.99.14 (3 matches)
Extended IP access list GUEST-OUT
    10 permit udp 192.168.40.0 0.0.0.255 host 192.168.20.10 eq domain (8841 matches)
    20 permit tcp 192.168.40.0 0.0.0.255 host 192.168.20.10 eq domain (12 matches)
    30 permit tcp 192.168.40.0 0.0.0.255 any eq www (40218 matches)
    40 permit tcp 192.168.40.0 0.0.0.255 any eq 443 (338104 matches)
    50 deny ip any any log (2277 matches)
\end{verify}

\begin{conceptbox}[The match counters are the diagnostic]
That output tells you more than the configuration does.

\begin{tight}
  \item \textbf{A rule with zero matches} is either unnecessary or unreachable
        --- shadowed by an earlier line. Both are worth investigating.
  \item \textbf{A rising counter on the final deny} tells you the ACL is dropping
        things, and \cmd{log} tells you what. This is the first place to look when
        an application ``sometimes'' fails.
  \item \textbf{Counters that do not move at all} mean the ACL is not in the path.
        Check it is actually applied, and in the direction you think.
\end{tight}

\cmd{clear access-list counters} before a test gives you a clean baseline.
\end{conceptbox}

\begin{verify}{R1}{show ip interface GigabitEthernet0/0/2 | include access list}
  Outgoing access list is not set
  Inbound  access list is GUEST-OUT
\end{verify}

\begin{pitfall}
\begin{tight}
  \item \textbf{Forgetting the implicit deny.} Permit what you need and everything
        else stops. This is the single biggest cause of ACL outages.
  \item \textbf{A general permit above a specific deny.} First match wins, so the
        deny never runs.
  \item \textbf{A standard ACL applied near the source.} Blocks the host from
        everything, not from the one destination you meant.
  \item \textbf{Port qualifier on the wrong side.} \cmd{eq 80} belongs with the
        destination for a client reaching a web server.
  \item \textbf{Writing the ACL and not applying it.} No error, no effect.
  \item \textbf{Applying it in the wrong direction.} \cmd{in} and \cmd{out} are
        relative to the interface, not to the building.
  \item \textbf{\cmd{no access-list 110} on a live numbered ACL.} Deletes the
        whole list; the interface is left denying everything.
  \item \textbf{Blocking your own management session.} Think about \cmd{vty}
        before you press enter on a remote device.
  \item \textbf{Forgetting supporting protocols.} DHCP, DNS and ICMP all need
        explicit permits in a restrictive list, or the VLAN looks completely
        dead.
\end{tight}
\end{pitfall}

\begin{breakfix}[An ACL meant to stop one subnet reaching the finance server has cut that subnet off from everything.]
The requirement was: hosts in \cmd{192.168.30.0/24} must not reach the finance
server at \cmd{192.168.20.15}. Everything else should be unaffected. After the
change, VLAN~30 cannot reach the internet, the file server, or DNS. It can still
reach hosts within its own VLAN.

\begin{verify}{R1}{show access-lists 20}
Standard IP access list 20
    10 deny   192.168.30.0, wildcard bits 0.0.0.255 (18442 matches)
    20 permit any (0 matches)
\end{verify}

\begin{verify}{R1}{show ip interface Vlan30 | include access list}
  Inbound  access list is 20
  Outgoing access list is not set
\end{verify}

\textbf{Diagnosis.} Two things are wrong and they compound.

First, this is a \emph{standard} ACL, so it can only match the source. It has no
way to express ``to the finance server'' --- the destination is invisible to it.

Second, it is applied \textbf{inbound on VLAN 30's own SVI}, which is as close to
the source as it is possible to get. Every packet leaving VLAN~30 for anywhere
hits line 10 and is dropped. Traffic \emph{within} the VLAN still works because
it is switched and never reaches the SVI at all --- which is exactly why the
symptom looks like a routing fault rather than an ACL.

The confirming detail is line 20: \cmd{permit any} has \textbf{zero matches}. Not
one packet from anywhere got past line 10. An ACL whose last line never matches
is dropping everything it sees.

\textbf{Fix.} The requirement names a destination, so it needs an extended ACL,
and an extended ACL can safely sit near the source.

\begin{config}{an ACL that expresses the actual requirement}
ip access-list extended VLAN30-POLICY
 10 remark Deny only the finance server, permit the rest
 20 deny   ip 192.168.30.0 0.0.0.255 host 192.168.20.15
 30 permit ip any any
!
interface Vlan30
 no ip access-group 20 in
 ip access-group VLAN30-POLICY in
\end{config}

Note line 30. Without it the extended ACL would reproduce the original outage by
a different route. Then \cmd{clear access-list counters} and re-test: line 20
should accumulate matches only when someone tries the finance server, and line 30
should carry everything else.
\end{breakfix}

\begin{lab}{Writing, reading and repairing ACLs}{Packet Tracer, or CML}
\textbf{Topology.} \cmd{R1} routing four VLANs --- staff 10, servers 20, lab 30,
guest 40 --- with a server at \cmd{192.168.20.10} running DNS and HTTP, a finance
server at \cmd{192.168.20.15}, and an internet-facing link.

\begin{tasks}
  \item Confirm full connectivity between all four VLANs before applying any ACL.
        Record the tests you used; you will repeat them after every change.
  \item Write and apply a numbered standard ACL restricting management access to
        the VTY lines to the staff VLAN only. Test from staff and from lab.
  \item Convert it to a named ACL. Insert an additional permitted host between
        two existing lines without retyping the list.
  \item Write an extended named ACL for the guest VLAN permitting only DNS to
        \cmd{192.168.20.10}, HTTP and HTTPS anywhere, and DHCP. Apply it inbound.
        Verify each permitted service works and at least three others do not.
  \item Read the match counters after your tests. Identify any rule with zero
        matches and explain why.
  \item \textbf{Break 1.} Remove the DHCP permit and release a guest host's
        lease. Predict the result before testing.
  \item \textbf{Break 2.} Move the \cmd{permit tcp ... eq 443} line above the
        \cmd{deny} that follows it, then below it, and observe the counters in
        each case.
  \item \textbf{Break 3.} Reproduce the Break \& Fix above deliberately: use a
        standard ACL near the source to try to protect one server. Confirm the
        symptom, then repair it with an extended ACL.
  \item \textbf{Break 4.} Apply an ACL that blocks your own SSH session on a
        device you can still reach by console. Recover it, then write down the
        habit that would have prevented it.
  \item \textbf{Extension.} Add \cmd{permit tcp any 192.168.20.0 0.0.0.255
        established} to an internet-facing inbound ACL. Verify that internal
        hosts can browse out and that inbound connections are refused, and
        explain why this is not a stateful firewall.
\end{tasks}
\end{lab}

\begin{checkpoint}
\begin{tightnum}
  \item Give the wildcard mask for \cmd{10.20.16.0/20} and the range it matches.
  \item Why is a standard ACL placed near the destination and an extended ACL
        near the source?
  \item An ACL's final \cmd{permit any} shows zero matches. What does that tell
        you?
  \item What does \cmd{established} match, and why is it not stateful?
  \item Name three supporting protocols that a restrictive inbound ACL must
        explicitly permit, and give the symptom for each if it is missed.
\end{tightnum}
\end{checkpoint}

\begin{chaptersummary}
\begin{tight}
  \item Top down, first match wins, implicit \cmd{deny any} at the end. Order is
        behaviour.
  \item Wildcard mask = 255 minus each octet of the subnet mask. Check it by
        adding it to the network address.
  \item Standard (1--99, 1300--1999) matches source only --- place near the
        destination. Extended (100--199, 2000--2699) matches source, destination,
        protocol and port --- place near the source.
  \item Named ACLs with sequence numbers can be edited a line at a time. Numbered
        ones on older IOS cannot, and deleting one that is still applied denies
        everything.
  \item Apply with \cmd{ip access-group in|out} on an interface, or
        \cmd{access-class in} on the VTY lines.
  \item \cmd{show access-lists} match counters are the real diagnostic: a zero on
        the last line means everything is being dropped; a zero in the middle
        means a shadowed rule.
  \item Remember DHCP, DNS and ICMP, or a restrictive ACL makes a VLAN look dead.
\end{tight}
\end{chaptersummary}

\begin{reviewq}
  \item \textbf{(MCQ)} Which wildcard mask matches exactly the range
        \cmd{172.20.32.0} to \cmd{172.20.39.255}?
        \begin{tight}
          \item[(a)] \cmd{0.0.7.255} \quad (b)~\cmd{0.0.15.255}
          \quad (c)~\cmd{0.0.3.255} \quad (d)~\cmd{255.255.248.0}
        \end{tight}
  \item \textbf{(MCQ)} A standard ACL should be applied:
        \begin{tight}
          \item[(a)] as close to the source as possible
          \item[(b)] as close to the destination as possible
          \item[(c)] on every interface of the router
          \item[(d)] inbound only
        \end{tight}
  \item \textbf{(Short)} Explain what happens when a numbered ACL that is applied
        to a live interface is deleted with \cmd{no access-list 110}.
  \item \textbf{(Short)} An ACL permitting HTTP and HTTPS is applied inbound on a
        user VLAN and the whole VLAN immediately loses connectivity, including
        the ability to obtain an address. Explain why, and give the two entries
        that were missing.
  \item \textbf{(Applied)} Write a named extended ACL that allows
        \cmd{10.10.10.0/24} to reach only SSH on \cmd{10.20.20.5}, DNS on
        \cmd{10.20.20.10}, and HTTPS anywhere, denying and logging everything
        else. State where you would apply it and in which direction, and justify
        both choices.
  \item \textbf{(Applied)} You are given an ACL and told users cannot print.
        Describe, in order, the four commands you would run to establish whether
        the ACL is responsible, and say what each one would prove.
\end{reviewq}
